\section{Projeto}
\label{sec:projeto}

\subsection{Função de transferência do somador inversor}
\label{sec:somador-prova}

Os blocos A e B são construídos com somadores inversores, portanto a
função de transferência dessa topologia é deduzida antes do projeto.
Considere um amplificador operacional ideal com realimentação negativa
por $R_f$, entradas $v_1,\dots,v_n$ ligadas ao nó inversor por
$R_1,\dots,R_n$ e entrada não inversora aterrada. Com ganho de malha
aberta infinito, a realimentação impõe $v^- = v^+ = 0$ (terra virtual)
e as correntes de entrada do amplificador são nulas.
A lei das correntes de Kirchhoff no nó inversor fornece
%
\begin{equation}
  \sum_{i=1}^{n}\frac{v_i - 0}{R_i} + \frac{v_{out} - 0}{R_f} = 0
  \quad\Longrightarrow\quad
  v_{out} = -R_f\sum_{i=1}^{n}\frac{v_i}{R_i}.
  \label{eq:somador}
\end{equation}
%
Cada entrada contribui com ganho $-R_f/R_i$, ajustável de forma
independente pela escolha de $R_i$. Coeficientes positivos são obtidos
com dois estágios em cascata, pois duas inversões restauram o sinal.

\subsection{Bloco A, somente somadores}

A função exigida,
$v_{out1} = -9v_{in1} + 0{,}2 - 3v_{in2} + 2{,}5v_{in3}$, tem dois
coeficientes negativos ($-9$ e $-3$), um positivo ($+2{,}5$) e um termo
constante de $+0{,}2$~V, e o roteiro restringe este bloco a somadores e
subtratores. A solução adotada foi uma cascata de dois somadores
inversores: $v_{in3}$ atravessa os dois estágios e sofre duas
inversões, chegando à saída com sinal positivo, enquanto $v_{in1}$ e
$v_{in2}$ entram apenas no segundo estágio e recebem sinal negativo.

O projeto seguiu as etapas abaixo, aplicando a
Eq.~\eqref{eq:somador} a cada estágio.

\begin{enumerate}
  \item \textbf{Primeiro estágio (XA1).} Ganho $-2{,}5$ para
        $v_{in3}$: com $R_{AF1}=25$~k e $R_{A4}=10$~k,
        $v_A = -(25/10)\,v_{in3} = -2{,}5\,v_{in3}$.
  \item \textbf{Segundo estágio (XA2).} Fixado
        $R_{AF2}=100$~k, os resistores de entrada resultam de
        $R_i = R_f/|c_i|$: $R_{A1} = 100/9 \approx 11{,}1111$~k para o
        coeficiente $-9$ e $R_{A2} = 100/3 \approx 33{,}3333$~k para o
        $-3$. A saída do primeiro estágio entra por
        $R_{A3}=100$~k (ganho $-1$), o que restaura o sinal positivo
        de $v_{in3}$.
  \item \textbf{Termo constante.} A constante $+0{,}2$~V foi
        implementada como uma entrada adicional do somador final ligada
        ao trilho de $-15$~V, já presente no circuito, evitando uma
        fonte adicional. O valor do resistor é obtido da própria
        Eq.~\eqref{eq:somador}: exige-se
        $-(R_f/R_{AR})(-15) = +0{,}2$~V, logo
        $R_{AR} = 15 R_f/0{,}2 = 75 R_f = 7{,}5$~M$\Omega$.
\end{enumerate}

A composição dos dois estágios comprova a função projetada:
%
\begin{align*}
  v_{out1} &= -\Bigl(\tfrac{100}{11{,}1111}v_{in1}
              + \tfrac{100}{33{,}3333}v_{in2}
              + \tfrac{100}{7500}(-15)
              + \tfrac{100}{100}v_A\Bigr)\\
           &= -\bigl(9v_{in1} + 3v_{in2} - 0{,}2 - 2{,}5v_{in3}\bigr)
            = -9v_{in1} + 0{,}2 - 3v_{in2} + 2{,}5v_{in3}.
\end{align*}
%
Os valores escolhidos estão na Tabela~\ref{tab:componentes} e
situam-se na faixa de dezenas de \si{\kilo\ohm}, suficientemente alta
para não carregar as fontes e suficientemente baixa para tornar
desprezíveis os efeitos das correntes de polarização. Como a
validação é por simulação,
foram usados os valores exatos ($100/9$ e $100/3$~\si{\kilo\ohm});
em uma montagem física esses valores seriam aproximados por associações
da série E96, como discutido na Seção~\ref{sec:discussao}.

\begin{figure}[htb]
  \centering
  \includegraphics[width=\textwidth]{rs2_bloco_a.png}
  \caption{Bloco A: somadores inversores XA1 e XA2, com as fontes de teste
  (senos de frequências diferentes).}
  \label{fig:esq-bloco-a}
\end{figure}

\subsection{Bloco B, topologia livre}

Para o bloco B a topologia é livre. Foi mantida a mesma cascata de
somadores inversores do bloco A, por uniformidade de projeto e porque a
estrutura de sinais é análoga: $v_{in4}$ e $v_{in5}$ exigem coeficiente
positivo (duas inversões) e $v_{in6}$ exige coeficiente negativo (uma
inversão). Aplicando a Eq.~\eqref{eq:somador} aos dois estágios:
%
\begin{align*}
  v_B &= -\Bigl(\tfrac{40}{20}v_{in4} + \tfrac{40}{10}v_{in5}\Bigr)
       = -(2v_{in4} + 4v_{in5}),\\
  v_{out2} &= -\Bigl(\tfrac{90}{10}v_{in6} + \tfrac{90}{90}v_B\Bigr)
       = 2v_{in4} + 4v_{in5} - 9v_{in6}.
\end{align*}

\begin{figure}[htb]
  \centering
  \includegraphics[width=\textwidth]{rs2_bloco_b.png}
  \caption{Bloco B: somadores inversores XB1 e XB2.}
  \label{fig:esq-bloco-b}
\end{figure}

\subsection{Bloco C, amplificador de instrumentação}
\label{sec:bloco-c}

O bloco C utiliza o amplificador de instrumentação clássico de três
amplificadores operacionais: dois
amplificadores de entrada (XC1, XC2) acoplados pelo resistor de ganho
$R_g$ e um subtrator de saída (XC3). A dedução do ganho usa as mesmas
hipóteses de amplificador ideal (curto virtual e correntes de entrada
nulas), com a notação de nós do esquemático
(Figura~\ref{fig:esq-bloco-c}).

\emph{Primeiro estágio.} O curto virtual impõe
$v_{ng1} = v_{out1}$ e $v_{ng2} = v_{out2}$, logo a corrente em $R_g$
vale $i = (v_{out1}-v_{out2})/R_g$. Como as entradas dos amplificadores
não drenam corrente, $i$ percorre também os dois resistores de
realimentação $R$ (RC1 e RC2):
%
\begin{align*}
  v_{na} &= v_{out1} + R\,i, \qquad v_{nb} = v_{out2} - R\,i,\\
  v_{na} - v_{nb} &= \Bigl(1 + \frac{2R}{R_g}\Bigr)(v_{out1}-v_{out2})
    = \Bigl(1 + \frac{2\cdot10\,\text{k}}{5\,\text{k}}\Bigr)
      (v_{out1}-v_{out2}) = 5\,(v_{out1}-v_{out2}).
\end{align*}

\emph{Segundo estágio.} No subtrator, o divisor $R_4$--$R_6$ fornece
$v_{np} = v_{nb}\,R_6/(R_4+R_6)$ e o curto virtual impõe
$v_{nn}=v_{np}$. A lei das correntes de Kirchhoff no nó inversor,
$(v_{na}-v_{nn})/R_3 = (v_{nn}-v_{out})/R_5$, resolve para
%
\[ v_{out} = \Bigl(1+\frac{R_5}{R_3}\Bigr)v_{nn}
             - \frac{R_5}{R_3}v_{na}
           = \frac{R_5}{R_3}\,(v_{nb} - v_{na})
           = 2{,}5\,(v_{nb}-v_{na}), \]
%
com $R_3=R_4=10$~k e $R_5=R_6=25$~k. Compondo os dois estágios:
%
\[ v_{out} = 2{,}5\cdot\bigl[-5(v_{out1}-v_{out2})\bigr]
           = 12{,}5\,(v_{out2}-v_{out1}), \]
%
que corresponde ao ganho especificado $G/2 = 25/2 = 12{,}5$. O ganho
foi fatorado em
$5\times2{,}5$ para que todos os resistores assumissem valores
comerciais exatos. Além do ganho, o amplificador de instrumentação
apresenta impedância de entrada elevada, não carregando os blocos A e
B, e rejeita sinais de modo comum, propriedade verificada na
Seção~\ref{sec:res-blocos}. A entrada $v_{out2}$ liga-se ao caminho não
inversor e $v_{out1}$ ao inversor, preservando o sinal de
$(v_{out2}-v_{out1})$.

\begin{figure}[htb]
  \centering
  \includegraphics[width=\textwidth]{rs2_bloco_c.png}
  \caption{Bloco C: amplificador de instrumentação (XC1, XC2, XC3) com
  as fontes do ensaio de ganho diferencial e modo comum.}
  \label{fig:esq-bloco-c}
\end{figure}

\subsection{Composição global}
\label{sec:composicao}

Substituindo as funções dos blocos A e B na do bloco C obtém-se a
função global implementada, que serve de referência para a validação:
%
\begin{align*}
  v_{out} &= 12{,}5\bigl[(2v_{in4}+4v_{in5}-9v_{in6})
             -(-9v_{in1}+0{,}2-3v_{in2}+2{,}5v_{in3})\bigr]\\
          &= 112{,}5v_{in1} + 37{,}5v_{in2} - 31{,}25v_{in3}
           + 25v_{in4} + 50v_{in5} - 112{,}5v_{in6} - 2{,}5.
\end{align*}

\begin{table}[htb]
  \centering
  \caption{Resistores dos três blocos.}
  \label{tab:componentes}
  \begin{tabular}{llll}
    \toprule
    Bloco & Referência & Valor & Papel / ganho \\
    \midrule
    A & RA4 / RAF1 & 10k / 25k & inversor de $v_{in3}$: $-2{,}5$ \\
    A & RA1 & 11{,}1111k & $v_{in1}$: $-100/11{,}1111 = -9$ \\
    A & RA2 & 33{,}3333k & $v_{in2}$: $-100/33{,}3333 = -3$ \\
    A & RAR & 7{,}5M & referência DC do trilho $-15$~V: $+0{,}2$~V \\
    A & RA3 / RAF2 & 100k / 100k & $v_A$: $-1$ \\
    \midrule
    B & RB1, RB2 / RBF1 & 20k, 10k / 40k & $v_{in4}$: $-2$; $v_{in5}$: $-4$ \\
    B & RB3, RB4 / RBF2 & 10k, 90k / 90k & $v_{in6}$: $-9$; $v_B$: $-1$ \\
    \midrule
    C & RC1, RC2 & 10k & realimentação do estágio de entrada \\
    C & RCG & 5k & $R_g$ (ganho de entrada 5) \\
    C & RC3, RC4 & 10k & entradas do subtrator \\
    C & RC5, RC6 & 25k & subtrator com ganho 2{,}5 \\
    \bottomrule
  \end{tabular}
\end{table}

\subsection{Modelo do LM741}
\label{sec:modelo}

O roteiro especifica o modelo LM741 Level 1 (modelo ideal) do
Micro-Cap. Como o LTspice não dispõe desse modelo, foi implementado um
subcircuito equivalente, \texttt{LM741\_LEVEL1}, com ganho de malha
aberta de $2\times10^5$, resistência de entrada de \SI{2}{\mega\ohm},
resistência de saída de \SI{75}{\ohm} e saída limitada pelas tensões de
alimentação, sem polos e sem limite de taxa de variação (\emph{slew
rate}), consistente com a definição de modelo ideal de nível 1. Os
parâmetros correspondem aos valores típicos da folha de dados do LM741
da Texas Instruments (documento SNOSC25D), reproduzida na
Figura~\ref{fig:datasheet-eletricas}.

\begin{verbatim}
.subckt LM741_LEVEL1 INP INN VCC VEE OUT params: Aol=200k Rin=2Meg Rout=75 Vhr=1
RIN INP INN {Rin}
BGAIN NDRV 0 V=limit({Aol}*(V(INP)-V(INN)), V(VEE)+{Vhr}, V(VCC)-{Vhr})
ROUT NDRV OUT {Rout}
COUT OUT 0 1p
.ends
\end{verbatim}

A alimentação de $\pm$\SI{15}{\volt}, inserida automaticamente pelo
Micro-Cap, foi incluída de forma explícita com as fontes VCC e VEE.
Para verificar se o modelo ideal não oculta efeitos relevantes, o
circuito completo também foi simulado com o macromodelo comercial do
LM741, que inclui polos, tensão de offset e taxa de variação de
\SI{0.5}{\volt\per\micro\second}. Como os estímulos do ensaio não
passam de \SI{1.5}{\kilo\hertz}, esses limites não são atingidos e
espera-se comportamento semelhante nos dois modelos.

\begin{figure}[htb]
  \centering
  \includegraphics[width=0.9\textwidth]{lm741-datasheet-eletricas.png}
  \caption{Características elétricas da folha de dados do LM741
  (Texas Instruments, SNOSC25D): ganho típico de
  \SI{200}{\volt\per\milli\volt}, resistência de entrada de
  \SI{2}{\mega\ohm} e excursão de saída de $\pm$\SI{14}{\volt} com
  alimentação de $\pm$\SI{15}{\volt}.}
  \label{fig:datasheet-eletricas}
\end{figure}

\subsection{Circuito completo}

A Figura~\ref{fig:esquematico} mostra o sistema com os sete LM741:
bloco A acima, bloco B abaixo e bloco C à direita, ligados pelos nós
\texttt{vout1} e \texttt{vout2}.

\begin{figure}[htb]
  \centering
  \includegraphics[width=\textwidth]{rs2_completo.png}
  \caption{Circuito completo (blocos A, B e C) com as fontes do teste
  global.}
  \label{fig:esquematico}
\end{figure}

\section{Procedimento de simulação}
\label{sec:simulacao}

As netlists e os esquemáticos são gerados pelo mesmo script Python
(Apêndice~\ref{ap:codigo}), o que garante que o circuito desenhado é o
mesmo que foi simulado. As simulações no domínio do tempo foram
organizadas em três grupos:

\begin{enumerate}
  \item \textbf{Resposta individual} (\texttt{rs2\_vin1.cir} a
        \texttt{rs2\_vin6.cir}): apenas uma entrada recebe um seno de
        \SI{50}{\milli\volt} e \SI{1}{\kilo\hertz}, com as demais
        aterradas. O seno tem pico positivo em
        $t=\SI{2.25}{\milli\second}$ e vale em
        \SI{2.75}{\milli\second}; a diferença entre os dois pontos,
        dividida por \SI{100}{\milli\volt}, fornece o ganho com sinal
        daquela entrada até a saída (e também até $v_{out1}$ e
        $v_{out2}$), eliminando o nível DC. A média da saída sobre
        ciclos inteiros mede o próprio nível DC, cujo valor teórico é
        $-12{,}5 \times 0{,}2 = -2{,}5$~V.
  \item \textbf{Função global} (\texttt{rs2\_completo.cir}): as seis
        entradas recebem senos de \SI{20}{\milli\volt} em frequências
        distintas (500 a \SI{1500}{\hertz}). Uma fonte comportamental
        calcula a função ideal em paralelo e mede-se o erro máximo e o
        erro RMS entre o circuito e a referência. As amplitudes foram
        escolhidas para que o pior caso ficasse em torno de
        \SI{7.3}{\volt}, abaixo dos \SI{14}{\volt} de excursão máxima,
        sem saturação.
  \item \textbf{Validação por bloco}
        (\texttt{rs2\_bloco\_a.cir}, \texttt{b}, \texttt{c}): cada
        bloco isolado recebe senos de frequências distintas e a saída é
        comparada com a respectiva função teórica. No bloco C
        aplica-se \SI{1}{\volt} de modo comum mais um diferencial
        senoidal de \SI{0.4}{\volt}, medindo-se o ganho diferencial e a
        rejeição de modo comum.
\end{enumerate}

O ensaio global foi repetido com o macromodelo comercial do LM741 e,
como verificação adicional, com um amplificador operacional moderno, o
ADA4610, para quantificar quanto do erro observado se deve às
limitações do LM741. As netlists correspondentes,
\texttt{rs2\_completo\_vendor.cir} e
\texttt{rs2\_completo\_moderno.cir}, e as demais estão reproduzidas
integralmente nos Apêndices~\ref{ap:netlist}
a~\ref{ap:variantes}.
